% THETA3-B29 contribution by /root/definite_bridge, 21 November 2026.
\subsection{THETA3-B29: actual theta sectors are orthogonal on the target}\label{B29}

Local Howe duality supplies an additional fact not used in B27:
the target images of its distinct source sectors are orthogonal.
Thus the inverse of the original fixed theta kernel is a finite
sum of positive-scalar multiples of algebraic sector inverses.
The scalar source periods remain.  This assertion concerns the
actual scalar oscillator coefficient of B17 and the actual
source sectors of B27, with the same spaces, splittings and
Haar measures throughout.

\paragraph{The local uniqueness theorem and the dual convention.}
Let $\sigma_\lambda$ denote the irreducible automorphic source
representation containing the coefficient functions in the
$\lambda$-sector of B27.  These are the representations of
$e_a(h)$, not of the functions integrated bilinearly against
them.  The latter lie in $\sigma_\lambda^\vee$.
For an irreducible target representation $\Pi$, the relevant
local occurrence is
\[
 \operatorname{Hom}_{G(F_v)\times S(F_v)}
             (\omega_v,\Pi_v\otimes\sigma_{\lambda,v})\ne0.
 \tag{B29a}
\]
In the convention where the local theta lift is the maximal
irreducible quotient of the other factor in
$\omega_v\twoheadrightarrow\sigma_{\lambda,v}\otimes
\Theta(\sigma_{\lambda,v})$, this says
$\Pi_v\simeq\theta(\sigma_{\lambda,v})$.
There is no contragredient on $\sigma_{\lambda,v}$ in
(B29a); its contragredient is the global bilinear input.

For a nonsplit finite place, Gan--Takeda,
\emph{A proof of the Howe duality conjecture}, Theorem 1.2,
printed pp.1--3, proves that the small theta lift is
irreducible or zero and that two nonzero small lifts agree
only when their inputs agree.  Its local field has
characteristic different from two and arbitrary residue
characteristic.  It therefore includes the ramified and
residue-characteristic-two completions needed here.
The spaces are Hermitian spaces over a quadratic field,
not quaternionic Hermitian spaces, so the separate
quaternionic qualification in that paper is irrelevant.

At a split finite place the groups form a type-II general
linear dual pair.  M\'inguez,
\emph{Correspondance de Howe explicite: paires duales de
type II}, Corollary 6.3(1), printed p.734, gives the
uniqueness of the other irreducible representation in
a nonzero quotient of the local oscillator representation.
Apply it also with the two members exchanged.  This is
the required injectivity for type II.  Fixed determinant
twists, the fixed inverse-transpose identification, and
the fixed nontrivial additive character of the actual
split model are those retained in B18.  They only apply
fixed bijections to the representation labels and do
not change uniqueness.

Only finite-place uniqueness will be needed below.
B27 already fixes the complete source infinity type,
so distinct source sectors differ at a finite place.
In particular the argument does not require a new
archimedean packet calculation or a separate global
irreducibility theorem for a theta lift.  It uses
irreducible quotients of the big local theta lift;
it never asserts that every big local theta lift is
itself irreducible.

\paragraph{A nonzero global constituent gives (B29a).}
The canonical cochain causes no additional representation
factor here.  B17 fixed a unit canonical frame $\kappa$
and the actual scalar oscillator vector
$\Phi=\Psi(\kappa)$, with
\[
   K(g,h)=\theta(\omega(g,h)\Phi).
 \tag{B29b}
\]
The archimedean coefficient pairing has already been
performed in this equality.  Applying a finite source
Hecke projector changes only the finite Schwartz input,
so $K_\lambda$ is still a vector in the same global
oscillator-to-automorphic theta map.  In particular it
is not merely a cohomology class compared to such a map.

Suppose a target irreducible constituent $\Pi$ occurs
in the automorphic representation generated by
$E_\lambda=\operatorname{im}T_\lambda$.
Equivalently the projection of the theta kernel to
the $\Pi$-isotypic target component and the
$\sigma_\lambda$-isotypic source component is nonzero.
Both quotients are compact, so all theta pairings used
to detect this projection converge absolutely.
Orthogonal automorphic projections commute with the
group actions.  Hence the theta map, followed by these
projections and by a choice of a nonzero irreducible
copy, gives a nonzero equivariant map from the global
oscillator representation to $\Pi\otimes\sigma_\lambda$.

Here is an explicit localization argument.  Fix a
finite place $v$.  Decompose the global Schwartz space
as the restricted tensor product with its $v$-factor.
A nonzero projected vector is detected on one finite
sum of pure tensors, hence on one choice of the
Schwartz vector away from $v$.  Pair the representation
factors away from $v$ with a suitable smooth dual
vector which detects a nonzero coefficient.  The
result is a nonzero $G(F_v)\times S(F_v)$-equivariant
map from $\omega_v$ to $\Pi_v\otimes
\sigma_{\lambda,v}$.  Admissibility ensures that a
vector fixed by an open compact subgroup lies in a
finite-dimensional fixed-vector tensor product, so
this is the usual smooth local Hom space.  This
proves (B29a) without a factorization formula for
the value of a global period.

The same convention can be checked directly by
integrating against $u\in\sigma_\lambda^\vee$ and
$v\in\Pi^\vee$:
\[
 \int_{[G]\times[S]}\theta_\Phi(g,h)u(h)v(g)\,dg\,dh.
 \tag{B29c}
\]
Changing variables after replacing $\Phi$ by
$\omega(g_0,h_0)\Phi$ moves $u,v$ by the inverse
right translations.  Thus this is an invariant
trilinear functional on
$\omega\otimes\sigma_\lambda^\vee\otimes\Pi^\vee$,
which is exactly (B29a), not its mislabeled dual.

\paragraph{Orthogonal target images and the actual full inverse.}
If one target constituent $\Pi$ occurred for both
$\lambda$ and $\eta$, local Howe uniqueness applied
to (B29a) would give
\[
             \sigma_{\lambda,v}\simeq
                   \sigma_{\eta,v}
                   \quad\hbox{for every finite }v.
 \tag{B29d}
\]
Their infinity types are already equal by B27, so the
global irreducible source representations are equal.
B27's multiplicity-one result then identifies their
finite Hecke sectors.  Consequently distinct sectors
cannot share a target irreducible automorphic
constituent.  The discrete unitary automorphic
decomposition is orthogonal, even if a target
representation has multiplicity.  It follows that
\[
 E_0=\mathop{\bigoplus}^{\perp}_{\lambda}E_\lambda,
       \qquad T_\lambda^*T_\eta=0
                       \quad(\lambda\ne\eta).
 \tag{B29e}
\]
Together with B27's source orthogonality this also
gives $T_\lambda T_\eta^*=0$.
Only nonzero sectors need be retained.  There are
at most $r=\dim E_0$ of them.

Use the matrices $A_\lambda$ of (B27i), with target
metric $G_B$.  Each is algebraic, positive
semidefinite and self-adjoint for $G_B$, and its
image is $E_\lambda$.  Equation (B29e) implies
\[
           A_\lambda A_\eta=0\quad(\lambda\ne\eta).
 \tag{B29f}
\]
Let $B_\lambda$ be the inverse of $A_\lambda$ on
$E_\lambda$, extended by zero on its $G_B$-orthogonal
complement.  This is an algebraic matrix: the image
and the complement are algebraic, and one inverts
an invertible algebraic matrix on the image.
Alternatively it is a polynomial in $A_\lambda$
which is zero at zero and is the reciprocal function
on its finitely many nonzero eigenvalues.
Put $P_\lambda=B_\lambda A_\lambda$; these are the
algebraic orthogonal projectors onto $E_\lambda$,
and their sum is the identity on $E_0$.

With every factor from B27 retained, the actual
inverse on the original fixed kernel image is
\[
 \boxed{\displaystyle
  Q_T^{-1}
    =\frac{1}{\nu_G|P_\chi|^2|\gamma|^2}
                 \sum_\lambda c_\lambda^{-1}B_\lambda,
          \qquad \gamma=\frac{C_x}{t_xt_y}.}
 \tag{B29g}
\]
Multiplication by
$Q_T=\nu_G|P_\chi|^2|\gamma|^2
\sum_\lambda c_\lambda A_\lambda$
gives $\sum_\lambda P_\lambda=1$, proving the formula.
Every $B_\lambda$ and $P_\lambda$ is realized by the
CM-source correspondence construction in B23, on
the preserved algebraic image supplied by B25.
This is an algebraic block decomposition of the
actual inverse.  It does not make the positive
numbers $c_\lambda$ algebraic or prove that their
ratios are algebraic.  B28 supplies a bound for
this inverse without requiring those ratios.

The projectors and all the scalars in (B29g) are
common to every source character or vector using
the same original kernel.  They have not been
chosen anew for each torus character.  Passing to
different complete types can change $P$ and the
kernel, and is outside this uniformity statement.

\paragraph{Exact transport along the original compact-place variation.}
There is a further uniformity statement for the
actual 032 parameter variation in B25.  Retain its
literal oscillator identity
\[
      \Phi_t=c_t\omega_S(s_t)\Phi_0,
                    \qquad c_t\ne0,
 \tag{B29h}
\]
where $s_t\in S(F\otimes\R)$ changes only the compact
target-place decompositions and all finite data and
the active canonical cochain stay fixed.  The scalar
$c_t$ includes the chosen Gaussian and constant
half-line normalizations.  It is retained exactly.
Then covariance of the scalar theta kernel and
right invariance of the original quotient Haar
measure give
\[
 \begin{aligned}
 K_t(g,h)&=c_tK_0(g,hs_t),\\
 T_t&=c_tT_0 R_S(s_t^{-1}),\\
 Q_{T_t}&=|c_t|^2Q_{T_0},
      &Q_{T_t}^{-1}&=|c_t|^{-2}Q_{T_0}^{-1}.
 \end{aligned}
 \tag{B29i}
\]
Here $(R_S(s)u)(h)=u(hs)$ is unitary on the full
source Hilbert space.  For the second line simply
substitute $h'=hs_t$ in the bilinear integral;
unitarity then proves the third line.  Thus even
the actual positive operator, not only its range,
is fixed up to the stated scalar along this
compact-place variation.  The finite source
projectors commute with $R_S(s_t)$, so (B29i)
holds in every sector separately and their target
images are unchanged.  No claim that the moved
source vector has the old compact type is needed;
the identity is on the full source Hilbert space.

Since the oscillator representation is unitary,
the modulus in (B29h) also satisfies
\[
                |c_t|^2=
                   \frac{\|\Phi_t\|_\omega^2}
                        {\|\Phi_0\|_\omega^2}.
 \tag{B29j}
\]
This does not discard a normalization: it identifies
it with the norms of the two specified Schwartz
vectors.  In a fixed real Schr\"odinger chart the
Gaussian identity
$\|\exp(\pi i x^{\mathsf t}Zx)\|_2^2
 =\det(2\operatorname{Im}Z)^{-1/2}$
computes each changing pure Gaussian factor.
Any originally included scalar half-line factors
must still be multiplied into this calculation.

\paragraph{An exact criterion for the remaining scalar gate.}
Retain only the nonzero sectors and fix one of them,
$\lambda_0$.  Write
$C=\nu_G|P_\chi|^2|\gamma|^2>0$ for the full
external factor in (B29g).  Then the following three
conditions are equivalent:
\[
 \begin{aligned}
 &Q_T=dM\quad\hbox{for some }d>0
                 \hbox{ and }M\in\operatorname{End}_{\overline\Q}(E_0);\\
 &Q_T^{-1}=d'M'\quad\hbox{for some }d'>0
                 \hbox{ and }M'\in\operatorname{End}_{\overline\Q}(E_0);\\
 &c_\lambda/c_{\lambda_0}\in
          \overline\Q\cap\R_{>0}\quad\hbox{for every retained }\lambda.
 \end{aligned}
 \tag{B29k}
\]
The first two conditions are equivalent by inversion.
If the third holds, put
$a_\lambda=c_\lambda/c_{\lambda_0}$; the exact formulas are
$Q_T=Cc_{\lambda_0}\sum_\lambda a_\lambda A_\lambda$
and
$Q_T^{-1}=(Cc_{\lambda_0})^{-1}
\sum_\lambda a_\lambda^{-1}B_\lambda$.
Their matrix factors are algebraic and their scalar
factors are positive, so no original normalization
has been removed.

Conversely, suppose $Q_T=dM$ as in the first condition.
Let $r_\lambda=\dim E_\lambda>0$.  Since
$B_\lambda A_\eta=0$ for $\lambda\ne\eta$ and
$B_\lambda A_\lambda=P_\lambda$, one has
\[
 \operatorname{Tr}(B_\lambda Q_T)=Cc_\lambda r_\lambda,
 \qquad
 \frac{c_\lambda}{c_{\lambda_0}}
   =\frac{r_{\lambda_0}}{r_\lambda}
      \frac{\operatorname{Tr}(B_\lambda M)}
           {\operatorname{Tr}(B_{\lambda_0}M)}.
 \tag{B29l}
\]
Every trace on the right is algebraic, and the
denominator is nonzero; the displayed first equality
also shows that their ratios are positive real.
This proves necessity.  In a single nonzero sector
the condition is vacuous.  With several sectors it
is precisely a period-ratio assertion, without any
remaining off-diagonal or multiplicity ambiguity.
It is not a consequence of the block decomposition
alone, and no general rationality conjecture is
being used to assert it.

\paragraph{Source scope.}
The new external inputs are the precise finite local
Howe uniqueness theorems above.  Gan--Takeda's
definitions, Theorem 1.2 and arbitrary-residue scope,
and the complete final proof in Section 6 were
read; M\'inguez's Theorem 5.1 with its complete
uniqueness proof and Corollary 6.3 were read.
The passage from the actual global kernel to local
Hom spaces, orthogonality, algebraic block inverse,
and variation identities are proved here.
No global irreducibility theorem, Galois
cohomological theta correspondence, or extra
\texttt{MainX} arithmetic family is inferred.
